\documentclass[12pt]{article}
\usepackage[utf8]{inputenc}
\usepackage[T1]{fontenc}
\usepackage{hyperref}
\title{Suruah\'a Language Resources for K-12 Curriculum Development}
\author{Compiled by Claude Sonnet 5.5 (claude-sonnet-5-5) for Andr\'as Kornai}
\date{}
\begin{document}
\maketitle

\section{Introduction}
Suruah\'a (also Zuruah\'a, Suruwah\'a; ISO 639-3 swx) is an Arawan language of Amazonas state, Brazil. Estimates of the speaker population are about 140 to 171 people. Search results describe the language as still largely undocumented.

\section{Primary Educational Resources}
No primers, textbooks or school materials were found.

\section{Linguistic and Reference Materials}
\begin{itemize}
\item Abstract for the University of Zurich colloquium talk on a grammatical description of Suruw\'ah\'a (Adriana Huber Azevedo's master's work): \url{https://www.comparativelinguistics.uzh.ch/en/colloquium/past-isle-colloquium/fall-term-2024/abstracts/2024-10-29.html}
\item UZH news item (2025) on the thesis prize: \url{https://www.comparativelinguistics.uzh.ch/en/news/2025/2025-04-21-herce.html}
\item Wikipedia entry with a preliminary phoneme inventory: \url{https://en.wikipedia.org/wiki/Zuruah%C3%A1_language}
\item Numeral system notes: \url{https://lingweb.eva.mpg.de/channumerals/Suruaha.htm}
\end{itemize}
No dictionary or published full grammar was found.

\section{Religious and Community Resources}
\begin{itemize}
\item Povos Ind\'igenas no Brasil (Instituto Socioambiental) page on the Zuruah\~a: \url{https://pib.socioambiental.org/en/Povo:Zuruahã}
\end{itemize}

\section{Digital and Audio Resources}
\begin{itemize}
\item Endangered Languages Project entry: \url{https://endangeredlanguages.com/node/85778}
\end{itemize}
No audio or video resources were found.

\section{Curriculum Development Context}
All community members are reported to use the language as a first language, but it is not known to be taught in schools. Written materials barely exist. Any curriculum would need to begin with community-led documentation and an agreed orthography, building on the new grammatical description. The community's own priorities and consent must come first.
\end{document}